% Final
\chapter{Seznam nefunkčních požadavků}\label{chapter:non-functional-requirements-list}

% Final
V této příloze lze nalézt seznam všech nefunkčních požadavků aplikace.


% Final
\section*{NP-1 Odstranění přebytečných funkcí}\label{sec:np-remove-extra-functions}

% Final
\begin{tabular}{lp{0.7\linewidth}}
    \textbf{Priorita:} & Must have \\
    \textbf{Kategorie:} & Udržovatelnost systému \\
\end{tabular}

% Final
Jelikož se nově vznikající aplikace bude zaměřovat na výuku programování, bude třeba odstranit funkce týkající se původní domény výuky jazyků. Tyto funkce jsem specifikoval v rámci analýzy stávajícího řešení.

% Final
\section*{NP-2 Nový design uživatelského rozhraní}\label{sec:np-new-ui-design}

% Final
\begin{tabular}{lp{0.7\linewidth}}
    \textbf{Priorita:} & Must have \\
    \textbf{Kategorie:} & Použitelnost systému \\
\end{tabular}

% Final
Jelikož aplikace mění své zaměření z výuky jazyků na výuku programování, je třeba kompletně změnit a přizpůsobit i celý vzhled uživatelského rozhraní. Měl by tak být vytvořen kompletně nový návrh uživatelského rozhraní, který se bude zahrnovat jak nové, tak i existující funkce aplikace.

% Final
\section*{NP-3 Zpřehlednění navigace v aplikaci}\label{sec:np-navigation-clarity}

% Final
\begin{tabular}{lp{0.7\linewidth}}
    \textbf{Priorita:} & Should have \\
    \textbf{Kategorie:} & Použitelnost systému \\
\end{tabular}

% Final
Aplikace by měla uživateli poskytovat přehlednější a jednodušší způsob navigace mezi jednotlivými stránkami. Z výsledků uživatelského testování v mé bakalářské práci \cite{Jenicek2024BachelorThesis}, která se zabývala původní aplikací, vyplynulo, že mají uživatelé problém s navigací mezi jednotlivými stránkami aplikace, zejména mezi jednotlivými sekcemi kurzů. Jelikož bude do aplikace také přidána řada nových funkcí, bude třeba navigaci aktualizovat tak, aby se uživatel mohl snadno orientovat v aplikaci a využívat všechny dostupné funkce.

% Final
\section*{NP-4 Vytvoření knihovny komponentů}\label{sec:np-component-library}

% Final
\begin{tabular}{lp{0.7\linewidth}}
    \textbf{Priorita:} & Should have \\
    \textbf{Kategorie:} & Rozšiřitelnost systému \\
\end{tabular}

% Final
Frontend aplikace využívá různé komponenty uživatelského rozhraní. Tyto komponenty by měly být přesunuty do samostatné knihovny tak, aby je bylo možné využít i v dalších částech projektu, jako například na webových stránkách nebo v systému pro správu obsahu aplikace.

% Final
Tato knihovna komponentů by měla stavět na již existujících komponentech a měla by využívat již zavedené technologie, tedy jazyk TypeScript a knihovny React, Storybook a Material UI.

% Final
\section*{NP-5 Zvýšení udržovatelnosti kódu}\label{sec:np-maintainability}

% Final
\begin{tabular}{lp{0.7\linewidth}}
    \textbf{Priorita:} & Should have \\
    \textbf{Kategorie:} & Udržovatelnost systému \\
\end{tabular}

% Final
Na základě analýzy stávající aplikace jsem se rozhodl přizpůsobit kód aplikace tak, aby byl lépe udržovatelný a rozšiřitelný. Konkrétní změny, které bude třeba provést, jsem popsal v rámci analýzy stávajícího řešení.

% Final
\section*{NP-6 Zavedení automatizovaných testů}\label{sec:np-automated-tests}

% Final
\begin{tabular}{lp{0.7\linewidth}}
    \textbf{Priorita:} & Should have \\
    \textbf{Kategorie:} & Použitelnost systému \\
\end{tabular}

% Final
Do frontendové části aplikace by měly být zavedeny automatizované testy s cílem zajištění správného fungování aplikace a zvýšení její kvality a použitelnosti.

% Final
Tento požadavek vychází jednak z analýzy existujícího řešení a jednak z návrhu úprav do budoucna, který jsem popsal v rámci mé bakalářské práce \cite{Jenicek2024BachelorThesis}, která se věnovala původní aplikaci.

% Final
\section*{NP-7 Automaticky generované třídy API}\label{sec:np-generated-api-classes}

% Final
\begin{tabular}{lp{0.7\linewidth}}
    \textbf{Priorita:} & Should have \\
    \textbf{Kategorie:} & Udržovatelnost systému \\
\end{tabular}

% Final
V kódu frontendové části aplikace by měly být automaticky generovány třídy reprezentující backendové API na základě OpenAPI specifikace \cite{OpenAPISpecification}, která je poskytována backendovou částí aplikace.

% Final
Požadavek vychází z analýzy existujícího řešení. Cílem tohoto požadavku je zjednodušit implementaci a zajistit konzistenci mezi frontendovou a backendovou částí aplikace.

% Final
\section*{NP-8 Animace a efekty}\label{sec:np-animations-effects}

% Final
\begin{tabular}{lp{0.7\linewidth}}
    \textbf{Priorita:} & Could have \\
    \textbf{Kategorie:} & Použitelnost systému \\
\end{tabular}

% Final
Do aplikace by měly být přidány základní animace a efekty s cílem zlepšit uživatelskou zkušenost a udržet pozornost uživatelů. Uživatelské rozhraní by tak mělo působit vizuálně poutavějším dojmem a mělo by výrazněji reagovat na interakce s uživatelem.

% Final
Tento požadavek vychází zejména z analýzy konkurenčních řešení, ve které konkurenční aplikace Mimo nebo Scrimba využívají různé animace a efekty.

% Final
\section*{NP-9 Nasazení aplikace jako PWA}\label{sec:np-pwa-deployment}

% Final
\begin{tabular}{lp{0.7\linewidth}}
    \textbf{Priorita:} & Should have \\
    \textbf{Kategorie:} & Použitelnost systému \\
\end{tabular}

% Final
Aplikace by měla být přizpůsobena tak, aby bylo možné ji sestavit a nasadit jako PWA, neboli jako progresivní webovou aplikaci. V důsledku toho by mělo být možné aplikaci stáhnout a nainstalovat na zařízení s operačními systémy Windows, Android či iOS. K aplikaci by tak mělo být možné přistupovat buď na webových stránkách, nebo přímo na zařízení uživatele.
